\chapter{Lateral Epicondylitis (Tennis Elbow)}
\index{Lateral Epicondylitis (Tennis Elbow)}
\label{sec:Lateral Epicondylitis (Tennis Elbow)}

\section{Overview}
Lateral epicondylitis, commonly known as tennis elbow, is a tendinopathy affecting the common extensor tendon origin at the lateral humeral epicondyle. It is one of the most common upper extremity overuse injuries, affecting 1--3\% of the general population annually. Despite its name, only 5--10\% of cases occur in tennis players; most are work-related or activity-related.
\section{Classification}

\begin{description}[font=\normalfont\bfseries,leftmargin=3em,labelwidth=2.5em]
  \item[Cause] Degenerative / Overuse
  \item[Organ System] Bones / Musculoskeletal
  \item[Pathophysiology] Angiofibroblastic tendinosis of the extensor carpi radialis brevis (ECRB) tendon origin
  \item[Duration] Chronic (weeks to months)
  \item[Onset] Gradual
  \item[Mode of Transmission] Non-communicable
  \item[Clinical Course] Chronic, self-limited in many cases
  \item[Severity] Mild to Moderate
  \item[Extent of Disease] Localized
  \item[Age Group] Adults (35--55 years)
  \item[Epidemiology] Incidence 4--7 per 1000 per year; equal gender distribution
  \item[ICD-11 Code] FA50.0
\end{description}

\section{Causes}
The condition results from repetitive microtrauma and overuse of the wrist extensor muscles, particularly the extensor carpi radialis brevis (ECRB). Histopathology reveals angiofibroblastic tendinosis with collagen disorganization, neovascularization, and mucoid degeneration rather than acute inflammation. Activities involving repetitive wrist extension, forearm pronation, or grip strengthening precipitate the condition.

\section{Risk Factors}

\begin{itemize}[noitemsep]
  \item Repetitive wrist extension activities (computer work, painting, carpentry)
  \item Racquet sports with poor technique or improper equipment
  \item Manual occupations (plumbers, mechanics, butchers)
  \item Age 35--55 years
  \item Smoking
  \item Obesity
  \item Previous upper extremity tendinopathy
\end{itemize}

\section{Signs and Symptoms}

\subsection{Early symptoms}
\begin{itemize}[noitemsep]
  \item Early manifestations of lateral epicondylitis (tennis elbow) may be subtle and nonspecific
\end{itemize}

\subsection{Common symptoms}
\begin{itemize}[noitemsep]
  \item \subsection{Common symptoms}
  \item \textbf{Common symptoms:}
  \item Pain over the lateral elbow, radiating into the forearm
  \item Pain aggravated by wrist extension against resistance (e.g., lifting, gripping, shaking hands)
  \item Tenderness directly over the lateral epicondyle
  \item Pain with passive wrist flexion (stretching the extensors)
  \item Weakness in grip strength
  \item Difficulty with simple tasks like holding a cup or turning a doorknob
  \item Pain or discomfort in the affected area
  \item Difficulty performing daily activities
\end{itemize}

\subsection{Less common symptoms}
\begin{itemize}[noitemsep]
  \item Atypical or less common presentations of lateral epicondylitis (tennis elbow) may occur
\end{itemize}

\subsection{Emergency warning signs}
\begin{itemize}[noitemsep]
  \item Severe or worsening symptoms of lateral epicondylitis (tennis elbow) require immediate medical evaluation
\end{itemize}
\section{Progression and Stages}
Tennis elbow typically develops gradually over weeks to months. In most cases, symptoms peak within 6--12 weeks and gradually resolve over 6--24 months with conservative management. A minority of patients develop chronic, persistent symptoms lasting beyond 12 months.

\section{Complications}

\begin{itemize}[noitemsep]
  \item Chronic pain and functional limitation
  \item Reduced grip strength and work capacity
  \item Tendon rupture (rare)
  \item Regional pain syndrome
  \item Reduced quality of life
  \item Emotional and psychological effects
\end{itemize}

\section{Diagnosis}

\begin{itemize}[noitemsep]
  \item Clinical diagnosis based on history and physical examination
  \item Positive Cozen test (pain with resisted wrist extension)
  \item Positive Mill test (pain with passive wrist flexion with elbow extended)
  \item Tenderness to palpation over the lateral epicondyle
  \item Ultrasound or MRI to confirm tendinosis and exclude other pathology
  \item X-ray to rule out intra-articular pathology
\end{itemize}

\section{Prevention}

\begin{itemize}[noitemsep]
  \item Use proper ergonomics and technique in sports and work activities
  \item Strengthen forearm extensor and flexor muscles
  \item Take frequent breaks during repetitive activities
  \item Use appropriate equipment (proper grip size, racquet string tension)
  \item Go for regular check-ups so any problems are caught early
\end{itemize}

\section{Treatment Options}

\begin{itemize}[noitemsep]
  \item Activity modification and relative rest from aggravating activities
  \item Physical therapy with eccentric strengthening exercises
  \item Nonsteroidal anti-inflammatory drugs (NSAIDs) for pain relief
  \item Counterforce bracing (tennis elbow strap) to reduce tendon strain
  \item Corticosteroid injections for short-term relief (controversial long-term)
  \item Platelet-rich plasma (PRP) injections (evidence for long-term benefit)
  \item Extracorporeal shockwave therapy for chronic cases
  \item Surgery (open or arthroscopic release) for refractory cases
\end{itemize}

\section{Living with It}

\begin{itemize}[noitemsep]
  \item Follow your treatment plan as prescribed
  \item Keep up with regular appointments with your healthcare provider
  \item Adjust your daily habits as needed
  \item Reach out to family, friends, or support groups for help
  \item Keep learning about your condition and how to manage it
\end{itemize}

\section{Outlook}
The prognosis for tennis elbow is favorable, with 80--90\% of patients improving with conservative management within 12--18 months. Corticosteroid injections provide rapid short-term relief but may increase recurrence risk. Surgical release is effective in 85--90\% of refractory cases. Prevention through proper ergonomics and conditioning is important, particularly for those in high-risk occupations.
